% Building 2, Part I. Words about words.
% One file for the whole part. Each chapter begins with a banner and \section*.

\part*{Part I. Words about words}
\phantomsection\addcontentsline{toc}{part}{Part I. Words about words}

Building 0, Chapter 0 writes rules about words; Building 1, Chapter 2 builds words and glues them. Here the words are statements and proofs, and the blurs are ``forget the proof'' and ``forget everything but the truth''.

%% ==================================================================
%%  Chapter 1  (proofs)
%% ==================================================================
\section*{Chapter 1. Proofs}
\phantomsection\addcontentsline{toc}{section}{\texorpdfstring{Chapter 1. Proofs \quad statements $\to$ what every world says}{Chapter 1. Proofs}}
\begin{center}
{\normalsize statements $\to$ what every world says \quad$\cdot$\quad from: Building 0, Chapters 0, 3; Building 1, Chapter 2 \quad$\cdot$\quad build, blur}\\[4pt]
{\small Statements and proofs are words. Glue two statements when each proves the other; glue them when every world gives them the same answer. The two blurs agree: that is completeness. For the single world of the whole numbers they do not: that is incompleteness.}
\end{center}

Written for statements built from letters with ``not'', ``and'', ``or'', ``if''. The lines for statements with ``for all'' are stated and reported (T4, T5); their proofs need a design of their own.

\subsection*{Given}

\begin{lines}
\item[\lab{G1}] (Building 0, Chapter 0, M3.) Words in a list of names; maps on words by recursion.
\item[\lab{G2}] (Building 1, Chapter 2.) Words, glued by rules from below; B2 by an action on words.
\item[\lab{G3}] (Building 0, Chapter 3.) $\{0, 1\}$ as $\Z/2$.
\end{lines}

\subsection*{Definitions}

\begin{lines}
\item[\lab{D1}] (Build, M3.) \emph{Statements}: words in letters $p_1, p_2, \dots$, the name $\bot$, and $\neg$, $\wedge$, $\vee$, $\to$. Write $A \leftrightarrow B$ for $(A \to B) \wedge (B \to A)$.
\item[\lab{D2}] (Build.) A \emph{proof} of $A$: a list of statements ending in $A$, each one an instance of $A \to (B \to A)$, \ $(A \to (B \to C)) \to ((A \to B) \to (A \to C))$, \ $(\neg B \to \neg A) \to (A \to B)$ (with $\neg A := A \to \bot$ and $\wedge$, $\vee$ written through $\to$, $\bot$), or got from two earlier ones $B$ and $B \to A$. Write $\vdash A$. Checking a proof is a recursion; finding one is a search.
\item[\lab{D3}] A \emph{world}: an arrow from the letters to $\{0, 1\}$. It gives every statement a value, by recursion ($\bot \mapsto 0$, and the usual tables).
\end{lines}

\subsection*{The blur}

On statements in the letters $p_1, \dots, p_n$. \ $O$: $\neg$, $\wedge$, $\vee$, $\to$.

\begin{lines}
\item[\lab{B1}] (From below: forget the proof.) $A \sim B$ when $\vdash A \leftrightarrow B$.
\item[\lab{B2}] (From above: forget all but the answers.) Questions: the value of a statement in each world.
\end{lines}

\subsection*{Propositions}

\begin{lines}
\item[\lab{P1}] (Survives.) $\neg$, $\wedge$, $\vee$, $\to$ all respect B1 and B2. The bags of B2 form $\{0, 1\}^{(\text{worlds})}$, with the operations taken world by world.
\item[\lab{P2}] (Dies: length; the proof.) $p$ and $\neg\neg p$ are glued, of lengths $1$ and $3$. ``Which proof'' is the first thing forgotten; ``how short a proof can be'' dies with it, and is a question of its own (R2).
\item[\lab{P3}] (Normal form.) Each bag of B2 holds exactly one ``or'' of full rows, one row for each world giving $1$ (in a fixed order). So there are $2^{2^n}$ bags.
\end{lines}

\subsection*{Theorems}

\begin{lines}
\item[\lab{T1}] (One way: sound.) If $A \sim B$ then $A$ and $B$ agree in every world. \ (Each axiom is $1$ in every world, and ``from $B$, $B \to A$ get $A$'' keeps ``$1$ in every world''.)
\item[\lab{T2}] (The other way. A construction.) If $A$ is $1$ in every world, $\vdash A$. \ (For each world $w$, prove ``the row of $w$ $\to$ $A$'' by recursion on $A$; then remove the letters one at a time with $\vdash (p \to A) \to ((\neg p \to A) \to A)$.)
\item[\lab{T3}] (Appears: completeness.) B1 and B2 agree. So the bags of B1 are exactly $2^{2^n}$, with P3 as normal forms, and ``provable'' is ``true in every world''.
\item[\lab{T4}] (Reported, classical: with ``for all''.) For statements about a floor with ``for all'' and ``there is'', and a list of axioms $\Gamma$: B1 (each proves the other from $\Gamma$) and B2 (same value in every floor obeying $\Gamma$) agree. The construction in the proof builds a floor out of the words themselves.
\item[\lab{T5}] (Reported, classical: Appears fails for one world.) Take the axioms of $\N$ (Building 0, Chapter 1, with $+$, $\times$) and replace B2 by the single world $\N$. Then B1 is strictly finer: some statement $G$ is true in $\N$ and is not glued to $\top$. Z4 holds, Z5 fails (Building 0, Chapter 0), with $G$ as witness. The proof codes words as numbers, so that statements about $\N$ speak about statements: a loop of the kind this building is named for.
\end{lines}

\subsection*{Remarks}

\begin{lines}
\item[\lab{R1}] (A twist.) A statement of the form ``for every $n$, a check on $n$ comes out yes'' can be refuted by one finite witness if it is false. So if it is shown to be neither provable nor refutable, it is true in $\N$. Many open questions about primes (Building 0, Chapter 18) have this form.
\item[\lab{R2}] (How much proof.) Two sizes, both what dies in B1: \emph{length}, the shortest proof, and \emph{strength}, the first rung of a ladder of axiom lists that proves $A$. The pigeonhole statement for $n + 1$ birds needs proofs of length growing faster than any power of $n$ in the weakest systems (classical, reported), and only a few lines once counting bags is allowed (Building 1, Chapter II.6 uses counts this way).
\item[\lab{R3}] (Old names.) D1: propositional formulas; D2: a Hilbert system; D3: valuations; B1: the Lindenbaum algebra, a Boolean algebra; T1: soundness; T2, T3: completeness (Post; Kalmár's proof); T4: Gödel's completeness theorem; T5: Gödel's first incompleteness theorem.
\item[\lab{R4}] (Machine check, \texttt{checks/b2-ch1-2-check.py}.) The axioms on many instances are $1$ in every world; $2^{2^2} = 16$ bags for two letters; P3 normal forms.
\end{lines}

\subsection*{Exercises}

\begin{lines}
\item[\lab{E1}] Write the normal form of P3 for $p \to q$ and for $\neg(p \leftrightarrow q)$.
\item[\lab{E2}] Show $\vdash p \to p$ from the axioms of D2 in five lines.
\item[\lab{E3}] Why does T2 not give a fast way to find proofs?
\end{lines}

\newpage

%% ==================================================================
%%  Chapter 2  (provability)
%% ==================================================================
\section*{Chapter 2. Equally provable}
\phantomsection\addcontentsline{toc}{section}{\texorpdfstring{Chapter 2. Equally provable \quad statements $\to$ how provable they are}{Chapter 2. Equally provable}}
\begin{center}
{\normalsize statements $\to$ how provable they are \quad$\cdot$\quad from: Chapter 1; Building 0, Chapters 18, 19, 34 \quad$\cdot$\quad blur, ladder}\\[4pt]
{\small Glue two statements when the theory proves ``this one is provable exactly when that one is''. Without letters the bags are a ladder of shape $\omega + 1$, and a statement's rung is the first depth at which it fails: a valuation, like $v_p$.}
\end{center}

All the lines about arithmetic here rest on three classical results, reported: the conditions on $\square$ (D1), L\"ob's theorem (G2), and the depth reading (G3). The chapter's own work is the blur and its reading.

\subsection*{Given}

\begin{lines}
\item[\lab{G1}] (Chapter 1, T5.) Statements of arithmetic can speak of proofs. For a statement $A$, $\square A$ is the statement ``$A$ is provable'', and $\vdash A$ gives $\vdash \square A$; $\vdash \square(A \to B) \to (\square A \to \square B)$; $\vdash \square A \to \square\square A$ (reported, classical).
\item[\lab{G2}] (Reported, classical: L\"ob.) $\vdash \square(\square A \to A) \to \square A$.
\item[\lab{G3}] (Reported, classical: depth.) For statements built from $\bot$ with $\neg$, $\wedge$, $\vee$, $\to$, $\square$ (no letters): read $A$ at each depth $n = 0, 1, 2, \dots$, with $\square A$ true at depth $n$ when $A$ is true at every depth below $n$. Then $\vdash A$ exactly when $A$ is true at every depth.
\end{lines}

\subsection*{The blur}

On statements of arithmetic. \ $O$: $\wedge$, $\square$.

\begin{lines}
\item[\lab{B1}] (From below.) $A \approx B$ when $\vdash \square A \leftrightarrow \square B$: ``equally provable''.
\item[\lab{B2}] (From above, without letters.) Question: $\rho(A) :=$ the first depth at which $A$ is false ($\infty$ if none), the \emph{rung} of $A$.
\end{lines}

\subsection*{Propositions}

\begin{lines}
\item[\lab{P1}] (Survives.) $\wedge$ and $\square$. \ ($\square(A \wedge B) \leftrightarrow \square A \wedge \square B$ by G1.)
\item[\lab{P2}] (Appears inside: L\"ob.) $A \approx (\square A \to A)$ for every $A$. \ (G1 one way, G2 the other.)
\item[\lab{P3}] (Dies: truth.) $\bot \approx \neg\square\bot$ (P2 with $A = \bot$). $\neg\square\bot$, ``consistency'', is true, and $\bot$ is false.
\item[\lab{P4}] (Dies: $\vee$ and $\neg$.) $\bot \approx \neg\square\bot$, but $\bot \vee \square\bot$ is not provable while $\neg\square\bot \vee \square\bot$ is.
\end{lines}

\subsection*{Theorems}

\begin{lines}
\item[\lab{T1}] (Appears, without letters: the bags are the rungs.) $A \approx B$ $\iff$ $\rho(A) = \rho(B)$. \ ($\square A$ is true at depth $n$ exactly when $n \le \rho(A)$; G3.)
\item[\lab{T2}] (A ladder of shape $\omega + 1$.) Every rung $k = 0, 1, 2, \dots$ and $\infty$ occurs: $\rho(\square^k \bot) = k$, and the provable statements are rung $\infty$. For every $A$ with $\rho(A) = k < \infty$: $\vdash \square A \leftrightarrow \square^{k+1}\bot$, so $\square^{k}\bot$ is a normal form of its bag. No two rungs are glued: that would give $\vdash \square^{k+2}\bot \to \square^{k+1}\bot$, and L\"ob would give $\vdash \square^{k+1}\bot$.
\item[\lab{T3}] (Survives, read on the rungs.) $\rho(\square A) = \rho(A) + 1$; \ $\rho(A \wedge B) = \min(\rho(A), \rho(B))$. \\
These are the laws of $v_p$ (Building 0, Chapter 18, D4) and of $\operatorname{ord}_a$ (Building 0, Chapter 34, P1) with ``$\times$'' read as ``$\wedge$'': a rung is the first place a statement stops being glued to the provable ones, as $v_p(a)$ is the first rung at which $a$ stops being glued to $0$.
\item[\lab{T4}] (Placement.) $\bot$ and consistency $\neg\square\bot$ on rung $0$; $\square\bot$ (``inconsistent'') on rung $1$; ``$\square\bot \to \square\square\bot$'' on rung $\infty$.
\end{lines}

\subsection*{Remarks}

\begin{lines}
\item[\lab{R1}] (What is known.) T1--T3 are the letterless normal form of the logic of $\square$ (Boolos, classical). Adding stronger boxes stretches the ladder to the ordinal $\varepsilon_0$, the measure of strength of arithmetic (Beklemishev, reported). The blur B1 with letters, and the mixture of B1 with ``same consequences of the form for every $n$'', are not worked out: open.
\item[\lab{R2}] (The twin.) The rung is a valuation: it is the fourth time the book meets ``first rung at which not glued'' (Building 0, Chapters 19, 21, 34). Here the ladder stops at $\omega + 1$ instead of having a bottom.
\item[\lab{R3}] (Old names.) $\square$: the provability predicate; G1: the Hilbert--Bernays--L\"ob conditions; G2: L\"ob's theorem; G3: Solovay's completeness for GL on letterless sentences, read on a line of worlds; P3: Gödel's second incompleteness theorem (consistency is not provable) seen as ``truth dies''.
\item[\lab{R4}] (Machine check, \texttt{checks/b2-ch1-2-check.py}.) On letterless statements up to size $3$ read at depths below $40$: T1 (bags are rungs), P2 (L\"ob instances), T3 ($\square$ climbs one rung, $\wedge$ is $\min$), T4, P4.
\end{lines}

\subsection*{Exercises}

\begin{lines}
\item[\lab{E1}] Find the rung of $\square\bot \to \bot$ and of $\square\square\bot \wedge \neg\square\bot$.
\item[\lab{E2}] Why is $\vee$ not $\max$ on rungs? Find $\rho(A \vee B)$ for $A = \bot$, $B = \square\bot$.
\item[\lab{E3}] Show from P2 that $A \approx \square A \to A \approx \square(\square A \to A) \to (\square A \to A)$.
\end{lines}
